\begin{problem}[理论考试][光纤光栅]{28}
  光纤光栅是一种介质折射率周期性变化的光学器件。设一光纤光栅的纤芯基体材料折射率为 $n_{1}=1.51$；在光纤中周期性地改变纤芯材料的折射率，其改变了的部分的材料折射率为 $n_{2}=1.55$；折射率分别为 $n_{2}$ 和 $n_{1}$、厚度分别为 $d_{2}$ 和 $d_{1}$ 的介质层相间排布，总层数为 N，其纵向剖面图如图(a)所示。在该器件设计过程中，一般只考虑每层界面的单次反射，忽略光在介质传播过程中的吸收损耗。假设入射光在真空中的波长为 $\lambda = 1.06\mu m$，当反射光相干叠加加强时，则每层的厚度 $d_{1}$ 和 $d_{2}$ 最小应分别为多少？若要求器件反射率达到 8\%，则总层数 N 至少为多少？

  提示：如图(b)所示，当光从折射率 $n_{1}$ 介质垂直入射到折射率 $n_{2}$ 介质时，界面上产生反射和透射，有：$\frac{\text{反射光电场强度}}{\text{入射光电场强度}} = \frac{n_1 - n_2}{n_1 + n_2}$，$\frac{\text{透射光电场强度}}{\text{入射光电场强度}} = \frac{2n_1}{n_1 + n_2}$，反射率 $= \left|\frac{\text{反射光电场强度}}{\text{入射光电场强度}}\right|^2$。
\end{problem}